\section{Semi-regular representations and physical parent kernels}
\label{sec:peps_semiregular_parent_equivalence}

Let $G$ be finite. A semi-regular representation contains every irreducible
representation of $G$. The comparison below derives the irreducible blocks
and their positive multiplicities from the original unitary representation,
using the representation decomposition and bond transformations of
\cite[Sections~4.1 and~7]{Schuch2010PEPS}.
The graph is finite and simple, its vertices are ordered, and each edge
$f=(t,h)$ has the resulting tail--head orientation.
Write $\mathcal G_R(A)=\operatorname{im}\mathcal C_R(A)$ for the range of
the open contraction with arbitrary virtual boundary coefficients, and
$\mathcal P_{\mathcal R}(A)$ for the common regional parent space.

\subsection{Irreducible coordinates and the regular representation}

\begin{theorem}[Positive irreducible multiplicities]%
    \label{thm:peps_semiregular_matrix_blocks}
    \lean{TNLean.PEPS.exists_unitary_matrix_blocks,
        TNLean.PEPS.exists_minimalSemiRegular_matrix}
    \leanok
    Let $U:G\to\operatorname{Mat}_X(\C)$ be unitary, with $X$ finite.
    There are a finite index set $I$, positive integers $d_i,m_i$,
    unitary irreducible representations $D_i$ of dimensions $d_i$ with
    pairwise distinct characters, and a matrix $Q$ such that
    \begin{align}
        Q^\dagger Q                       & =I,
                                          & U(g)  & =Q\left(\bigoplus_iD_i(g)\otimes I_{m_i}\right)Q^\dagger,
        \label{eq:peps_semiregular_matrix_blocks}                                                             \\
        \langle\chi_U,\chi_{D_i}\rangle_G & =m_i.
    \end{align}
    If $U$ is semi-regular, then $D=\bigoplus_iD_i$ is unitary and
    semi-regular. The dimensions, multiplicities, and coordinate matrix
    are conclusions, rather than additional assumptions on $U$.
\end{theorem}
\begin{proof}\leanok
    \uses{thm:peps_unitary_character_blocks,
        thm:peps_semiregular_remove_copies}
    Choose the orthonormal irreducible rows of
    Theorem~\ref{thm:peps_unitary_character_blocks}. The matrices of the
    action on one row of each component give $D_i$; the common action on
    its $m_i$ rows gives $D_i\otimes I_{m_i}$ after exchanging the row
    and matrix indices. The change between these orthonormal coordinates
    and the original basis is $Q$, giving
    (\ref{eq:peps_semiregular_matrix_blocks}). The character multiplicity
    equals the number of rows. Apply
    Theorem~\ref{thm:peps_semiregular_remove_copies} for semi-regularity;
    the direct sum of the unitary $D_i$ is unitary.
\end{proof}

\begin{theorem}[Removing copies preserves semi-regularity]%
    \label{thm:peps_semiregular_remove_copies}
    \lean{TNLean.PEPS.isSemiRegular_blockMatrixRepresentation_of_intertwiner}
    \leanok
    Suppose $U$ is semi-regular and
    $U(g)=Q(\bigoplus_iD_i(g)\otimes I_{m_i})Q^\dagger$.
    Then $\bigoplus_iD_i$ is semi-regular. This implication does not
    require $Q$ to be isometric, or the $D_i$ to be irreducible.
\end{theorem}
\begin{proof}\leanok
    \uses{thm:peps_semiregular_group_algebra,lem:peps_twirled_representation}
    Semi-regularity is equivalent to linear independence of the group
    operators. If $\sum_gc_gD(g)=0$, each diagonal block gives
    $\sum_gc_gD_i(g)=0$. Tensoring with $I_{m_i}$ and conjugating by $Q$
    yields $\sum_gc_gU(g)=0$. Hence every $c_g$ is zero.
\end{proof}

\begin{theorem}[The character after dimension multiplicities]%
    \label{thm:peps_semiregular_regular_character}
    \lean{TNLean.PEPS.character_blockMatrixRepresentation,
        TNLean.PEPS.irreducibleCharacterFinset_blockMatrixRepresentation,
        TNLean.PEPS.character_multiplicityRestoredRepresentation,
        TNLean.PEPS.character_multiplicityRestoredRepresentation_eq_regular,
        TNLean.PEPS.card_multiplicityRestored_eq}
    \leanok
    For any finite family of representations $D_i$ of dimensions $d_i$,
    put $D=\bigoplus_iD_i$ and $D^{[d]}=\bigoplus_iD_i\otimes I_{d_i}$.
    Then $\chi_D=\sum_i\chi_{D_i}$ and
    $\chi_{D^{[d]}}=\sum_id_i\chi_{D_i}$.
    If each $D_i$ is irreducible and $D$ is semi-regular, the occurring
    irreducible characters are exactly $\{\chi_{D_i}:i\in I\}$.
    If these characters are also pairwise distinct, then
    \begin{align}
        \chi_{D^{[d]}}(g) & =\chi_{\mathrm{reg}}(g)
        =\begin{cases}|G|,&g=1,\\0,&g\ne1,\end{cases}
                          & \sum_i d_i^2            & =|G|.
        \label{eq:peps_semiregular_regular_character}
    \end{align}
\end{theorem}
\begin{proof}\leanok
    \uses{thm:peps_semiregular_trace}
    Traces add under direct sums and multiply under tensor products.
    Character orthogonality shows that an irreducible character absent
    from all the summands has zero multiplicity. Conversely, semi-regularity
    includes every irreducible character of the regular representation.
    Its character expansion is therefore $\sum_id_i\chi_{D_i}$, with
    no repetitions. Evaluating at the identity gives the dimension formula.
\end{proof}

\begin{definition}[Fourier coefficient matrices]%
    \label{def:peps_semiregular_block_fourier}
    \lean{TNLean.PEPS.irreducibleBlockFourierMatrix,
        TNLean.PEPS.normalizedIrreducibleBlockFourierMatrix}
    \leanok
    For $a,b\in\{0,\ldots,d_i-1\}$ and $g\in G$, define
    \begin{align}
        T_{(i,a,b),g} & =D_i(g)_{a,b},
                      & \widehat T_{(i,a,b),g} & =sqrt{\frac{d_i}{|G|}}D_i(g)_{a,b}.
        \label{eq:peps_semiregular_block_fourier}
    \end{align}
\end{definition}

\begin{theorem}[Fourier identification with the regular representation]%
    \label{thm:peps_semiregular_block_fourier}
    \lean{TNLean.PEPS.irreducibleBlockFourierMatrix_injective,
        TNLean.PEPS.irreducibleBlockFourierMatrix_intertwining,
        TNLean.PEPS.nonempty_equiv_multiplicityRestored_leftRegular,
        TNLean.PEPS.exists_invertible_multiplicityRestored_leftRegular,
        TNLean.PEPS.exists_isometry_multiplicityRestored_leftRegular}
    \leanok
    \uses{def:peps_semiregular_block_fourier}
    For arbitrary $D_i$, $TL(g)=D^{[d]}(g)T$, where $L$ is the left
    regular representation. If $D$ is semi-regular, $T$ is injective.
    If the $D_i$ are also irreducible with pairwise distinct characters,
    $T$ is invertible and $D^{[d]}$ and $L$ are equivalent: there are
    matrices $S,R$ with $RS=SR=I$ and $L(g)=SD^{[d]}(g)R$.
    If the $D_i$ are unitary, one may take $Q=\widehat T^\dagger$, giving
    \begin{align}
        Q^\dagger Q=I,\qquad QQ^\dagger=I,\qquad
        L(g)=QD^{[d]}(g)Q^\dagger.
        \label{eq:peps_semiregular_unitary_fourier}
    \end{align}
\end{theorem}
\begin{proof}\leanok
    \uses{thm:peps_semiregular_regular_character,lem:peps_twirled_representation}
    Multiplicativity of $D_i$ gives $TL(g)=D^{[d]}(g)T$.
    If $Tc=0$, all matrix coefficients of $\sum_gc_gD(g)$ vanish;
    semi-regularity forces $c=0$. The dimension identity in
    (\ref{eq:peps_semiregular_regular_character}) makes $T$ bijective.
    For unitary blocks, direct multiplication gives
    \begin{align}
        (\widehat T^\dagger\widehat T)_{g,h}
        =|G|^{-1}\sum_id_i\operatorname{tr}D_i(g^{-1}h)
        =\delta_{g,h}.
    \end{align}
    Equal dimensions give the other inverse identity. The row weights
    commute with $D^{[d]}$, so the intertwining identity also holds for
    $\widehat T$, proving~(\ref{eq:peps_semiregular_unitary_fourier}).
\end{proof}

\subsection{The oriented incident action}

\begin{definition}[Incident representations and their numbered coordinates]%
    \label{def:peps_semiregular_incident_action}
    \lean{TNLean.PEPS.graphIncidentMatrix,
        TNLean.PEPS.graphIncidentRepresentation,
        TNLean.PEPS.graphNumberedVirtualConfigEquiv,
        TNLean.PEPS.graphNumberedVirtualCoeffEquiv,
        TNLean.PEPS.graphNumberedIncidentRepresentation,
        TNLean.PEPS.graphNumberedPhysicalCoeffEquiv}
    \leanok
    Let $U:G\to\operatorname{Mat}_X(\C)$ be a representation.
    On the incident coordinate space $\mathcal V_v=\C^{X^{\operatorname{Inc}(v)}}$,
    put
    \begin{align}
        \rho_v(g)_{\sigma,\eta}
        =\prod_{f\in\operatorname{Inc}(v)}
        \begin{cases}
            U(g^{-1})_{\eta_f,\sigma_f}, & v=t(f), \\
            U(g)_{\sigma_f,\eta_f},      & v=h(f).
        \end{cases}
        \label{eq:peps_semiregular_incident_action}
    \end{align}
    Thus a head carries $U(g)$ and a tail carries $U(g^{-1})^{\mathsf T}$.
    A numbering $\{0,\ldots,|X|-1\}\cong X$ gives a coefficient-space
    isomorphism $J_v$ from numbered virtual labels to $\mathcal V_v$.
    Define the numbered action by $\rho_v^{\mathrm{num}}=J_v^{-1}\rho_vJ_v$.
    A physical enumeration $e_v:X^{\operatorname{Inc}(v)}\cong
        \{0,\ldots,p-1\}$ similarly gives $E_v:\mathcal V_v\cong\C^p$.
\end{definition}

\begin{theorem}[The canonical averaging tensor is G-injective]%
    \label{thm:peps_semiregular_averaging_g_injective}
    \lean{TNLean.PEPS.graphIncidentMatrix_one,
        TNLean.PEPS.graphIncidentMatrix_mul,
        TNLean.PEPS.graphIncidentRepresentation_apply,
        TNLean.PEPS.graphNumberedIncidentRepresentation_apply,
        TNLean.PEPS.regularSiteMap_graphAveragingSite,
        TNLean.PEPS.isGInjective_graphAveragingSite,
        TNLean.PEPS.localTensorMap_numberedGraphDressedAveragingTensor_one,
        TNLean.PEPS.isGInjective_numberedGraphDressedAveragingTensor_one}
    \leanok
    \uses{def:peps_semiregular_incident_action,def:peps_g_injective}
    The matrices in~(\ref{eq:peps_semiregular_incident_action}) satisfy
    $\rho_v(1)=I$ and $\rho_v(gh)=\rho_v(g)\rho_v(h)$, and act on
    coefficient functions by matrix multiplication.
    The unweighted averaging site has local map
    $P_v=|G|^{-1}\sum_g\rho_v(g)$.
    It is $G$-injective for $\rho_v$.
    In numbered coordinates its local map is $A_v=E_vP_vJ_v$ and is
    $G$-injective for $\rho_v^{\mathrm{num}}$.
    No unitarity or semi-regularity assumption is required here.
\end{theorem}
\begin{proof}\leanok
    Products of independent incident indices separate into products of
    matrix products. At a tail, inversion and transposition each reverse
    the order, giving multiplicativity. Expanding the averaging site
    yields $P_v$. Averaging gives $P_v\rho_v(g)=P_v$ and fixes every
    invariant vector, hence is injective on the invariant subspace.
    Conjugation by $J_v$ transports invariance, and $E_v$ is injective;
    therefore $E_vP_vJ_v$ has the same two $G$-injectivity properties.
\end{proof}

\subsection{Equivalence of the complete parent spaces}

\begin{definition}[Common incident enumerations]%
    \label{def:peps_semiregular_incident_enumeration}
    \lean{TNLean.PEPS.countedIncidentEnumeration}
    \leanok
    If $|\operatorname{Inc}(v)|=k$ at every vertex and $X$ is finite,
    choose the finite-set enumeration
    $e^X_v:X^{\operatorname{Inc}(v)}\cong\{0,\ldots,|X|^k-1\}$.
    Write $A(U,e)$ for the unweighted canonical averaging tensor in a
    supplied enumeration $e$, and put $A_{\mathrm{reg}}=A(L,e^G)$.
\end{definition}

\begin{theorem}[Arbitrary semi-regular canonical parent spaces]%
    \label{thm:peps_semiregular_canonical_parent_equivalence}
    \lean{TNLean.PEPS.nonempty_canonicalSemiRegularParentEquiv,
        TNLean.PEPS.finrank_regionParentGroundSpace_arbitrarySemiRegular_eq,
        TNLean.PEPS.nonempty_ker_regionParentHamiltonian_arbitrarySemiRegular_equiv}
    \leanok
    \uses{def:peps_semiregular_incident_enumeration}
    Suppose one unitary semi-regular representation $U$ is used on all
    edges, $|\operatorname{Inc}(v)|=k$ for every vertex, and every edge
    has both endpoints in some region $R_i$. For any common physical
    enumeration $e$ of the incident $X$-labels, there is a linear
    isomorphism
    \begin{align}
        \mathcal P_{\mathcal R}(A(U,e))
        \cong\mathcal P_{\mathcal R}(A_{\mathrm{reg}}).
        \label{eq:peps_semiregular_canonical_parent_equivalence}
    \end{align}
    The dimensions are equal. For a finite region family, the kernels
    of any positive local parent Hamiltonians with these exact regional
    ranges are also linearly isomorphic.
\end{theorem}
\begin{proof}\leanok
    \uses{thm:peps_semiregular_matrix_blocks,
        thm:peps_semiregular_block_fourier,
        thm:peps_copy_positive_parent,
        thm:peps_fourier_coordinate_parent,thm:peps_region_parent_common_kernel}
    Derive $D_i,d_i,m_i,Q$ from
    Theorem~\ref{thm:peps_semiregular_matrix_blocks}.
    Transport canonical parents from $U$ to $D^{[m]}$ by $Q^\dagger$.
    Theorem~\ref{thm:peps_copy_positive_parent} changes the positive
    copy family $m_i$ to $d_i$. The normalized Fourier matrix then
    transports $D^{[d]}$ to $L$ by
    Theorem~\ref{thm:peps_fourier_coordinate_parent}.
    Constant incidence supplies the intermediate common alphabets, and
    edge coverage supplies the supported inverse for the copy change.
    Compose these three isomorphisms. Positivity identifies the kernel
    of each parent Hamiltonian with its common regional parent space.
\end{proof}

\begin{theorem}[Physical G-injective tensors and their canonical parents]%
    \label{thm:peps_semiregular_physical_canonical_parent}
    \lean{TNLean.PEPS.nonempty_gInjectiveCanonicalParentEquiv}
    \leanok
    \uses{def:peps_semiregular_incident_action,def:peps_g_injective}
    Let $A_v:\mathcal V_v^{\mathrm{num}}\to\C^p$ be $G$-injective
    for $\rho_v^{\mathrm{num}}$, and let $e$ enumerate the canonical
    physical incident space, possibly with another dimension $q$.
    If every vertex belongs to a region $R_i$, then
    $\mathcal P_{\mathcal R}(A)\cong
        \mathcal P_{\mathcal R}(A(U,e))$.
    Neither constant incidence nor semi-regularity is needed for this
    fixed-representation comparison, provided $e$ is supplied.
\end{theorem}
\begin{proof}\leanok
    \uses{thm:peps_semiregular_averaging_g_injective,
        thm:peps_range_parent_g_comparison,
        thm:peps_range_parent_equivalence}
    The canonical averaging tensor is $G$-injective for the same
    incident action. The fixed-representation comparison supplies
    physical maps $F_v$ that are injective on its site images and
    satisfy $A_v=F_vA(U,e)_v$. Vertex coverage makes the restriction of
    $\bigotimes_vF_v$ an isomorphism of the parent spaces. Take its inverse.
\end{proof}

\begin{theorem}[Complete kernels for semi-regular G-injective tensors]%
    \label{thm:peps_semiregular_g_injective_parent_kernel}
    \lean{TNLean.PEPS.nonempty_semiRegularGInjectiveParentEquiv,
        TNLean.PEPS.finrank_regionParentGroundSpace_semiRegularGInjective_eq,
        TNLean.PEPS.nonempty_ker_regionParentHamiltonian_semiRegularGInjective_equiv}
    \leanok
    \uses{def:peps_semiregular_incident_action,
        def:peps_semiregular_incident_enumeration,def:peps_region_parent_interaction}
    Let $U:G\to\operatorname{Mat}_X(\C)$ be one unitary semi-regular
    representation used uniformly on the edges of the finite ordered
    simple graph. Suppose $|\operatorname{Inc}(v)|=k$ at every vertex
    and $A_v:\mathcal V_v^{\mathrm{num}}\to\C^p$ is $G$-injective
    for the actual oriented action $\rho_v^{\mathrm{num}}$.
    Assume that every vertex is in some $R_i$, and that every edge has
    both endpoints in some $R_i$. Then
    \begin{align}
        \mathcal P_{\mathcal R}(A) & \cong
        \mathcal P_{\mathcal R}(A_{\mathrm{reg}}),
                                   & \dim\mathcal P_{\mathcal R}(A) & =
        \dim\mathcal P_{\mathcal R}(A_{\mathrm{reg}}).
        \label{eq:peps_semiregular_physical_parent_equivalence}
    \end{align}
    If the region family is finite, let $H_i,K_i\ge0$ satisfy
    $\ker H_i=\mathcal G_{R_i}(A)$ and
    $\ker K_i=\mathcal G_{R_i}(A_{\mathrm{reg}})$.
    Their sums over the corresponding region placements obey
    $\ker H\cong\ker K$.
    The physical dimension $p$ need not equal $|G|^k$.
\end{theorem}
\begin{proof}\leanok
    \uses{thm:peps_semiregular_physical_canonical_parent,
        thm:peps_semiregular_canonical_parent_equivalence,thm:peps_region_parent_common_kernel}
    Use the common incident enumeration $e^X$ and compose
    \begin{align}
        \mathcal P_{\mathcal R}(A)
        \cong\mathcal P_{\mathcal R}(A(U,e^X))
        \cong\mathcal P_{\mathcal R}(A_{\mathrm{reg}}).
    \end{align}
    The first comparison uses vertex coverage; the second uses edge
    coverage, constant incidence, and semi-regularity.
    Dimensions are preserved by linear isomorphisms.
    For positive interactions, $\langle\psi,H\psi\rangle=0$ holds
    exactly when every regional term annihilates $\psi$; hence
    $\ker H=\mathcal P_{\mathcal R}(A)$ and likewise for $K$.
\end{proof}

% Scope: these ordered-simple-graph, uniform-U statements do not themselves
% identify the native right/up torus closures in SCP10 Theorems 5.7/5.9.
% The source allows link-dependent representations (source lines 1310--1316).
% Reversing an edge replaces U by its contragredient; arbitrary multiplicities
% need not be dual-symmetric. Native plaquette coverage, orientation matching,
% and identification with the named commuting-pair closure vectors are separate
% requirements. Small-period tori may identify parallel bonds in a simple graph.
% No spectral-gap or excited-spectrum equivalence is asserted here.
